%%%
%
% $IR-25-01-ToF $
% $Datum: 10.6.25 $
% $Pfad: ArduinoLab $
% $Dateiname: ArduinoLab.Tex
% $Version: 4621 $
%
% !TeX spellcheck = de_DE/GB
% !TeX program = pdflatex
% !BIB program = biber/bibtex
% !TeX encoding = utf8
%
%%%


\chapter{Arduino Lab for MicroPython} \label{ArduinoLab}

\section{Einführung}

Arduino Lab for MicroPython ist ein Editor für MicroPython-Programme. Arduino Lab for MicroPython unterstützt die Verbindung mit einem Board, das Hochladen von Programmen, den Dateitransfer und eine interaktive REPL-Shell.

\section{Installation}

\begin{enumerate}
	\item\textbf{Stellen Sie den Ein/Aus-Schalter des Alvik Roboters auf „Aus“.} 
	\item\textbf{Verbinden Sie den ESP32 auf dem Alvik über ein USB-C-Kabel mit dem Computer.} 
	\item\textbf{Laden Sie das Arduino Lab for MicroPython von der offiziellen Arduino Website herunter.} 
	\item\textbf{Greifen Sie über „File Manager“ auf den Ordnerinhalt zu oder wechselt zu einem neuen Ordner.Wie in der Abbildung\ref{Filemanager} dargestellt.} 	
	\begin{figure}[ht]  
		\begin{center}
			\includegraphics[width=12cm]{Arduino/Alvik/Filemanager}
			\caption{Arduino lab for MicroPython File manager} 
			\label{Filemanager}                                                                                                                                                                                                                                                                                                          
		\end{center}
	\end{figure}
	\item\textbf{Klicken Sie auf die Schaltfläche „Connect“ und wählen Sie den richtigen COM-Anschluss.Wie in der Abbildung\ref{Arduinolabconnect} dargestellt.} 
	\begin{figure}[ht]  
		\begin{center}
			\includegraphics[width=12cm]{Arduino/Alvik/Arduinolabconnect}
			\caption{Arduino lab for MicroPython Connect} 
			\label{Arduinolabconnect}                                                                                                                                                                                                                                                                                                          
		\end{center}
	\end{figure}
	\item\textbf{Klicken Sie auf die Schaltflächen „Editor und REPL“. Der obere wei\ss e Teil ist der Editor zum Schreiben von Skripten und der untere schwarze Teil ist das Terminal.Wie in der Abbildung\ref{EditorandREPL} dargestellt.} 
	\begin{figure}[ht]  
		\begin{center}
			\includegraphics[width=12cm]{Arduino/Alvik/EditorAndREPL}
			\caption{Editor und REPL} 
			\label{EditorandREPL}                                                                                                                                                                                                                                                                                                          
		\end{center}
	\end{figure}
	\item\textbf{Klicken Sie auf die Schaltfläche Run, um das Standard-Beispielskript auszuführen. Nachdem das Skript ausgeführt wurde, wird``Hallo, MicroPython!'' auf dem Terminal ausgegeben.Wie in der Abbildung\ref{Run} dargestellt.} 
	\begin{figure}[ht]  
		\begin{center}
			\includegraphics[width=12cm]{Arduino/Alvik/Run}
			\caption{Arduino lab for MicroPython Run} 
			\label{Run}                                                                                                                                                                                                                                                                                                          
		\end{center}
	\end{figure}
	\item\textbf{Öffnen Sie den „File Manager“ eine neue Datei. Benennen und öffnen Sie die neue Datei.} 
	\item\textbf{Schieben Sie den Ein/Aus-Schalter des Alvik in die Position Ein.} 
	
\end{enumerate}


\section{Erste Schritte}
Der Alvik kann zunächst programmiert werden, indem Sie die Text ein- und ausgegeben werden.Das Programm~\ref{Eingabetext} ist in der Abbildung dargestellt.
\begin{figure}[ht]
	\begin{center}
		\includegraphics[width=12cm]{Arduino/Alvik/Input}
		\caption{Eingabe und Ausgabe von Text Programm} 
		\label{Eingabetext}                                                                                                                                                                                                                                                                                                          
	\end{center}
\end{figure}

\begin{enumerate}
	\item \textbf{Eingabe und Ausgabe von Text}
	\begin{itemize}
		\item \textbf{Toolbox verwenden:}Importieren von sleep aus dem Zeitmodul und Verwendung der Funktion \PYTHON{sleep()}, um das Programm einige Sekunden warten zu lassen, bevor es mit der nächsten Programmzeile fortfährt.
		\item \textbf{Eingabe- und Ausgabezeichen:}Die Funktion \PYTHON{input()} wird zur Eingabe von Informationen in das Terminal verwendet. Alle vom Benutzer eingegebenen Informationen werden in der Variablen student\_name gespeichert. Hier wartet das Programm auf die Eingabe und fährt erst dann fort, wenn die Eingabe eingegangen ist. Verwenden \PYTHON{print()}, um Text auf dem Terminal auszugeben.
		\item \textbf{Verkettung:}Verwenden den + Symbol, um die Zeichenkette mit der Variable student\_name zu kombinieren.
		\item \textbf{Einstellung der Interaktionsgeschwindigkeit:}Verwenden von \PYTHON{sleep()}, um die Interaktion zu verlangsamen, was zu einem natürlicheren Dialogfluss führt.
		
	\end{itemize}
\end{enumerate}	

\section{Datei im Arduino Alvik speichern}

Der Arduino Alvik ermöglicht das Ausführen integrierter Programme, indem Dateien vom Computer über Arduino Lab für MicroPython in den Arduino Alvik importiert werden.
\begin{itemize}
	\item Öffnen Sie die Arduino Lab for MicroPython-Datei über die Schaltfläche „File Manager“. Wie in der Abbildung\ref{alab} gezeigt, werden im linken Fenster die in Arduino Alvik integrierten Dateien und im rechten Fenster die auf dem Computer gespeicherten Dateien angezeigt.
	\begin{figure}[ht]
		\begin{center}
			\includegraphics[width=12cm]{Arduino/Alvik/alab}
			\caption{Arduino Lab for MicroPython Dateispeicherort} 
			\label{alab}                                                                                                                                                                                                                                                                                                          
		\end{center}
	\end{figure}
	\item Wählen Sie ein Programm aus und klicken Sie auf die Schaltfläche ``Hochladen'' (\textless)\ref{hochladen}, um das Programm von dem Computer auf das Arduino Alvik hochzuladen.
	\begin{figure}[ht]
		\begin{center}
			\includegraphics[width=12cm]{Arduino/Alvik/Upload}
			\caption{Schaltfläche ``Hochladen''} 
			\label{hochladen}                                                                                                                                                                                                                                                                                                          
		\end{center}
	\end{figure}
	\item Öffnen Sie die in Arduino Alvik integrierte Datei main.py \ref{main}, löschen Sie den darin enthaltenen Code, geben Sie dem Programm \PYTHON{import xxx} ein (wobei xxx der Name der zu importierenden Datei ist) und speichern anschlie\ss end.
	\begin{Python} [caption={Code für main.py},label={main}]
		import xxx #(wobei xxx der Name der zu importierenden Datei ist)
	\end{Python}
	\item Trennen das Arduino Alvik vom Computer und schalten es ein. Das Arduino Alvik führt automatisch die zuvor eingerichteten eingebauten Programme aus.
	\item Beenden das Programm durch Ausschalten des Schalters am Arduino Alvik.
\end{itemize}



